\documentclass[10pt,a4paper,twocolumn]{article}

% ─── Encoding & fonts ────────────────────────────────────────────
\usepackage[utf8]{inputenc}
\usepackage[T1]{fontenc}
\usepackage{lmodern}
\usepackage{microtype}

% ─── Language ────────────────────────────────────────────────────
\usepackage[spanish,english]{babel}

% ─── Page geometry ───────────────────────────────────────────────
\usepackage[top=2.5cm, bottom=2.5cm, left=2cm, right=2cm,
            columnsep=0.6cm]{geometry}

% ─── Colour palette ──────────────────────────────────────────────
\usepackage[dvipsnames,table]{xcolor}
\definecolor{ecogreen}{RGB}{22,101,52}
\definecolor{ecolightgreen}{RGB}{187,247,208}
\definecolor{ecoteal}{RGB}{15,118,110}
\definecolor{ecogray}{RGB}{75,85,99}
\definecolor{ecobg}{RGB}{240,253,244}
\definecolor{ecored}{RGB}{185,28,28}
\definecolor{ecoyellow}{RGB}{180,120,0}
\definecolor{codeblue}{RGB}{37,99,235}
\definecolor{codebg}{RGB}{248,250,252}

% ─── Graphics & diagrams ─────────────────────────────────────────
\usepackage{graphicx}
\usepackage{tikz}
\usetikzlibrary{shapes.geometric, arrows.meta, positioning, fit,
                backgrounds, shadows.blur, calc, decorations.markings}

% ─── Tables ──────────────────────────────────────────────────────
\usepackage{booktabs}
\usepackage{tabularx}
\usepackage{multirow}
\usepackage{array}

% ─── Code listings ───────────────────────────────────────────────
\usepackage{listings}
\lstset{
  basicstyle=\ttfamily\footnotesize,
  keywordstyle=\color{codeblue}\bfseries,
  commentstyle=\color{ecogray}\itshape,
  stringstyle=\color{ecored},
  backgroundcolor=\color{codebg},
  frame=single,
  rulecolor=\color{ecogray!40},
  breaklines=true,
  numbers=left,
  numberstyle=\tiny\color{ecogray},
  xleftmargin=1.2em,
  framexleftmargin=1.2em,
  captionpos=b
}

% ─── Hyperlinks ──────────────────────────────────────────────────
\usepackage[colorlinks=true,
            linkcolor=ecogreen,
            citecolor=ecoteal,
            urlcolor=ecoteal,
            pdftitle={EcoMarket: A Microservices-Based E-Commerce Platform},
            pdfauthor={Josep Vladimir Garate Quispe}]{hyperref}
\usepackage{url}

% ─── Section styling ─────────────────────────────────────────────
\usepackage{titlesec}
\titleformat{\section}{\large\bfseries\color{ecogreen}}
  {\thesection.}{0.5em}{}[\vspace{-2pt}\color{ecogreen}\rule{\linewidth}{0.6pt}]
\titleformat{\subsection}{\normalsize\bfseries\color{ecoteal}}
  {\thesubsection}{0.5em}{}
\titleformat{\subsubsection}{\small\bfseries\color{ecogray}}
  {\thesubsubsection}{0.5em}{}

% ─── Abstract box ────────────────────────────────────────────────
\usepackage{abstract}
\renewcommand{\abstractnamefont}{\bfseries\color{ecogreen}}
\renewcommand{\abstracttextfont}{\small\itshape}

% ─── Float control ───────────────────────────────────────────────
\usepackage{float}
\usepackage{caption}
\captionsetup{font=small, labelfont={bf,color=ecoteal}, skip=4pt}

% ─── Misc ────────────────────────────────────────────────────────
\usepackage{enumitem}
\usepackage{amsmath}
\usepackage{mdframed}
\usepackage{multicol}

% ─── Header / footer ─────────────────────────────────────────────
\usepackage{fancyhdr}
\pagestyle{fancy}
\fancyhf{}
\renewcommand{\headrulewidth}{0.4pt}
\renewcommand{\headrule}{\hbox to\headwidth{\color{ecogreen}\leaders\hrule height \headrulewidth\hfill}}
\fancyhead[L]{\small\textcolor{ecogray}{EcoMarket — Taller de Software, UNA-Puno}}
\fancyhead[R]{\small\textcolor{ecogray}{\thepage}}

% ═══════════════════════════════════════════════════════════════════
%  TITLE
% ═══════════════════════════════════════════════════════════════════
\title{%
  \vspace{-1cm}%
  {\color{ecogreen}\large\bfseries\scshape SoftwareX — Academic Report}\\[4pt]
  \rule{\linewidth}{1.5pt}\\[8pt]
  {\LARGE\bfseries EcoMarket: A Microservices-Based\\[4pt]
   E-Commerce Platform for Ecological Products\\[4pt]
   with AI-Driven Chemical Auditing}\\[6pt]
  \rule{\linewidth}{1.5pt}
}

\author{%
  \textbf{Josep Vladimir Garate Quispe}\\
  \small Facultad de Ingeniería Estadística e Informática\\
  \small Universidad Nacional del Altiplano, Puno, Perú\\[4pt]
  \small \href{mailto:josep.garate@est.unap.edu.pe}{\texttt{josep.garate@est.unap.edu.pe}}\\[6pt]
  \small\textbf{Course:} Taller de Software \quad
  \small\textbf{Instructor:} Torres Cruz Fred
}

\date{\small June 2026}

% ═══════════════════════════════════════════════════════════════════
%  DOCUMENT
% ═══════════════════════════════════════════════════════════════════
\begin{document}

\twocolumn[
  \maketitle
  \vspace{-0.3cm}
  \begin{@twocolumnfalse}
    \begin{abstract}
      EcoMarket is a comprehensive e-commerce platform tailored for the Puno
      region of Peru. It connects local ecological and organic producers directly
      with environmentally conscious consumers through a scalable microservices
      architecture. The platform integrates an AI-based chemical auditing engine
      that analyzes product ingredients, assigns a quantitative Eco-Score
      (0--100), generates color-coded ecological badges, and produces
      tamper-proof PDF certificates anchored with SHA-256 hashes on an
      immutable MongoDB ledger. The system is built with Next.js 16 (PWA
      frontend), Java Spring Boot, Python FastAPI, and Node.js Express as
      backend services, PostgreSQL and MongoDB as databases, Kong as API
      gateway, and is containerized with Docker. Native Peruvian payment
      methods (Yape, Plin, TuPay) and international Stripe payments are
      supported. The platform is deployed on Vercel (frontend) and Render
      (backend services) and is publicly accessible at
      \href{https://eco-market-final-software-s2w6.vercel.app}{\texttt{eco-market-final-software-s2w6.vercel.app}}.
      Source code is available at
      \href{https://github.com/JosepVl123463/EcoMarketFinalSoftware}{\texttt{github.com/JosepVl123463/EcoMarketFinalSoftware}}
      under the MIT License.
    \end{abstract}
    \vspace{0.3cm}
    \noindent\textbf{Keywords:}
    e-commerce; microservices; artificial intelligence; chemical auditing;
    Next.js; Spring Boot; FastAPI; ecological products; Puno; Perú
    \vspace{0.5cm}
    \hrule
    \vspace{0.4cm}
  \end{@twocolumnfalse}
]

% ───────────────────────────────────────────────────────────────────
\section{Motivation and Significance}
% ───────────────────────────────────────────────────────────────────

\subsection{Context and Problem Statement}

The high-altitude region of Puno, Peru, sustains rich agricultural
biodiversity — native quinoa, cañihua, native potatoes, and a wide
range of organic products. Despite this potential, producers face two
persistent barriers: (1) limited access to digital markets and heavy
reliance on intermediaries that compress profit margins, and (2) widespread
consumer distrust caused by unverifiable ``ecological'' product claims.

EcoMarket was created to address both barriers simultaneously by providing
a vertically integrated platform: direct-to-consumer digital commerce
combined with algorithmic ecological certification.

\subsection{Gap in Existing Solutions}

General-purpose e-commerce platforms (Mercado Libre, Amazon) lack
region-specific ecological focus, local payment method integration,
and automated chemical verification. EcoMarket fills this gap by
combining commerce with verifiable AI-based auditing in a single system.

\subsection{Contribution}

The key contributions are: (i) a production-grade microservices
architecture tailored for Peruvian ecological commerce; (ii) an
AI chemical auditing engine with a curated database of 11 prohibited
chemical families; (iii) tamper-proof audit certificates via SHA-256
immutable hashes; and (iv) integration of local Peruvian payment methods
alongside international Stripe processing.

% ───────────────────────────────────────────────────────────────────
\section{Software Description}
% ───────────────────────────────────────────────────────────────────

\subsection{System Architecture Overview}

EcoMarket follows a domain-driven microservices design. Six backend
services communicate through a central API Gateway; each service owns
its data store and exposes a REST API. Figure~\ref{fig:arch} depicts
the full architecture.

% ─── TikZ architecture diagram ─────────────────────────────────────
\begin{figure*}[ht]
\centering
\begin{tikzpicture}[
  node distance=0.55cm and 0.7cm,
  font=\small,
  box/.style={
    rectangle, rounded corners=4pt, draw, thick,
    minimum width=2.3cm, minimum height=0.75cm,
    text centered, align=center
  },
  frontend/.style={box, fill=ecolightgreen!70, draw=ecogreen!80},
  gateway/.style={box, fill=ecoteal!20, draw=ecoteal, minimum width=4.5cm},
  service/.style={box, fill=blue!10, draw=codeblue, minimum width=2.2cm},
  db/.style={
    cylinder, shape border rotate=90, draw, thick,
    fill=yellow!15, draw=ecoyellow!80,
    minimum width=1.4cm, minimum height=0.65cm,
    text centered, font=\scriptsize
  },
  arrow/.style={-{Latex[length=2.5mm]}, thick, color=ecogray},
  darrow/.style={<->, thick, dashed, color=ecogray!60},
]

% ── Row 0: Frontend ──────────────────────────────────────────
\node[frontend, minimum width=7cm] (fe)
  {\textbf{Frontend — Next.js 16 PWA (Vercel)}\\
   \scriptsize Catalog · Auth · Checkout · Admin · Producer Dashboard};

% ── Row 1: Gateway ───────────────────────────────────────────
\node[gateway, below=0.9cm of fe] (gw)
  {\textbf{API Gateway — Kong / Node.js (:8000)}\\
   \scriptsize Rate-Limiting · CORS · JWT Routing · Load Balancing};

% ── Row 2: Microservices ──────────────────────────────────────
\node[service, below left=1.1cm and 2.8cm of gw] (auth)
  {\textbf{Auth}\\
   \scriptsize Spring Boot\\
   \scriptsize :8081};

\node[service, below left=1.1cm and 0.7cm of gw] (product)
  {\textbf{Product}\\
   \scriptsize Spring Boot\\
   \scriptsize :8082};

\node[service, below right=1.1cm and 0.7cm of gw] (payment)
  {\textbf{Payment}\\
   \scriptsize Node.js\\
   \scriptsize :8083};

\node[service, below right=1.1cm and 2.8cm of gw] (audit)
  {\textbf{Audit}\\
   \scriptsize FastAPI\\
   \scriptsize :8084};

% ── Row 3: AI + Notification ─────────────────────────────────
\node[service, below=1.0cm of audit, fill=purple!10, draw=purple!60] (ai)
  {\textbf{AI Engine}\\
   \scriptsize FastAPI\\
   \scriptsize :8085};

\node[service, below=1.0cm of payment, fill=orange!10, draw=orange!70] (notif)
  {\textbf{Notification}\\
   \scriptsize FastAPI\\
   \scriptsize :8086};

% ── Row 4: Databases ─────────────────────────────────────────
\node[db, below=0.9cm of auth] (pg1)
  {PostgreSQL\\Users};

\node[db, below=0.9cm of product] (pg2)
  {PostgreSQL\\Products};

\node[db, below=0.9cm of audit, fill=green!10, draw=ecogreen!70] (mongo)
  {MongoDB\\Audits};

\node[db, below=0.9cm of ai, fill=red!8, draw=red!50] (redis)
  {Redis\\Cache};

% ── Arrows ───────────────────────────────────────────────────
\draw[arrow] (fe) -- (gw);

\draw[arrow] (gw) -- (auth);
\draw[arrow] (gw) -- (product);
\draw[arrow] (gw) -- (payment);
\draw[arrow] (gw) -- (audit);

\draw[arrow] (payment) -- (notif);
\draw[arrow] (audit) -- (ai);
\draw[darrow] (payment) -- (product);
\draw[darrow] (audit) -- (product);

\draw[arrow] (auth) -- (pg1);
\draw[arrow] (product) -- (pg2);
\draw[arrow] (audit) -- (mongo);
\draw[arrow] (ai) -- (redis);

\end{tikzpicture}
\caption{EcoMarket microservices architecture. Six domain services communicate
  through a centralized API Gateway. Arrows indicate primary data-flow direction;
  dashed bidirectional arrows denote inter-service REST calls.}
\label{fig:arch}
\end{figure*}

\subsection{Frontend — Next.js 16 PWA}

The frontend is a Progressive Web Application built with
Next.js~16, React~19, TypeScript~5, and Tailwind CSS~4.
State is managed by Zustand~5; server-state is handled by
React Query~5 with Axios for HTTP communication.
Key screens include: product catalog with Eco-Score filters,
producer registration (multi-step), admin audit dashboard,
and a multi-method checkout. Security features include
Cloudflare Turnstile CAPTCHA, XSS prevention via DOMPurify,
JWT Bearer token injection via Axios interceptors, and
strict CSP headers.

\subsection{Backend Microservices}

\subsubsection{Auth Service — Spring Boot (:8081)}

Handles user registration (customer and producer multi-step flows),
login, JWT issuance, and role-based access control.
Key entities: \texttt{User} (id, email, role, ecoScore) and
\texttt{Provider} (id, userId, ruc, businessName, certification).
Tokens carry a 24-hour expiration, encoded with HS256.
Rate limiting is enforced at 30~req/min and 200~req/hour per IP.

\subsubsection{Product Service — Spring Boot (:8082)}

Manages the product catalog, inventory, and orders.
Core entities: \texttt{Product} (UUID PK, category, status,
eco\_score CHECK~0--100, images TEXT[]) and \texttt{Order}
with multi-vendor \texttt{OrderItem} support.
Products follow a status lifecycle: \texttt{PENDING}
$\rightarrow$ \texttt{APPROVED} / \texttt{REJECTED}.
A 15\,\% platform fee is recorded in the
\texttt{payment\_splits} table for revenue distribution.

\subsubsection{Payment Service — Node.js/Express (:8083)}

Orchestrates all payment flows. For international cards, it creates
a Stripe Checkout Session and verifies webhook signatures. For local
methods, it simulates Yape, Plin, and TuPay flows, then calls
product-service to confirm the order and notification-service to
push a push notification. Rate limiting is set at 15~payments/min.

\subsubsection{Audit Service — FastAPI/Python (:8084)}

The audit service is the technical centrepiece of EcoMarket.
It receives product ingredient lists, runs them against an internal
chemical database of 11 restricted families, computes an Eco-Score,
assigns ecological badges, and stores an immutable audit record in
MongoDB. It also generates downloadable PDF certificates
using ReportLab.

\subsubsection{AI Engine — FastAPI/Python (:8085)}

Provides a mock computer-vision and NLP interface for ingredient
detection from product images. The architecture is designed to be
swapped with a production model (e.g., a fine-tuned
CLIP or BERT-based classifier) without changing the inter-service
contract.

\subsubsection{Notification Service — FastAPI/Python (:8086)}

A thin dispatcher that formats and delivers push notifications
(Firebase FCM) on payment confirmation and audit result events.

% ───────────────────────────────────────────────────────────────────
\section{AI Chemical Auditing Engine}
% ───────────────────────────────────────────────────────────────────

\subsection{Chemical Database}

The engine encapsulates a curated database of 11 chemical families
considered incompatible with ecological certification, each with
aliases, risk level, reason for concern, and regulatory reference
(Table~\ref{tab:chemicals}).

\begin{table}[H]
\centering
\caption{Excerpt of the EcoMarket chemical restriction database.}
\label{tab:chemicals}
\renewcommand{\arraystretch}{1.1}
\scriptsize
\begin{tabularx}{\columnwidth}{>{\bfseries}l l l}
\toprule
Family & Risk & Regulation \\
\midrule
\rowcolor{red!8}   Sulfatos (SLS/SLES) & High & CE~1223/2009 Ann.~III \\
\rowcolor{red!8}   Parabenos           & High & CE~1223/2009 Ann.~V  \\
\rowcolor{red!8}   Ftalatos            & High & REACH Reg.~1907/2006  \\
\rowcolor{red!8}   BPA                 & High & CE~10/2011            \\
\rowcolor{orange!8}Petrolatos          & Med. & COSMOS Standard       \\
\rowcolor{orange!8}Siliconas           & Med. & ECOCERT Protocol      \\
\rowcolor{orange!8}Triclosan           & Med. & CE~1223/2009 Ann.~VI  \\
\rowcolor{orange!8}Formaldehído        & Med. & CE~1223/2009 Ann.~III \\
\rowcolor{orange!8}Microplásticos      & Med. & EU Plastic Reg.       \\
\rowcolor{orange!8}Colorantes Azoicos  & Med. & CE~1223/2009 Ann.~IV  \\
\rowcolor{orange!8}Aluminio            & Med. & EFSA 2008 Report      \\
\bottomrule
\end{tabularx}
\end{table}

\subsection{Scoring Algorithm}

The Eco-Score $S$ is computed as:
\begin{equation}
  S = \max\!\left(0,\; 100 - \sum_{i} p_i\right)
\end{equation}
where $p_i = 25$ for each high-risk ingredient detected and
$p_i = 10$ for each medium-risk ingredient.
A product is \textbf{APPROVED} if $S \geq 70$; otherwise
\textbf{REJECTED}.

\subsection{Badge Assignment Logic}

Figure~\ref{fig:badges} shows the badge assignment thresholds.

\begin{figure}[H]
\centering
\begin{tikzpicture}[font=\scriptsize, node distance=0.4cm]
  \tikzset{
    badge/.style={rectangle, rounded corners=3pt, draw, thick,
                  minimum width=2.6cm, minimum height=0.5cm,
                  text centered, fill=ecolightgreen!60, draw=ecogreen}
  }
  \node[badge] (b1) {\textbf{Eco-Friendly} \quad $S \geq 90$};
  \node[badge, below=of b1] (b2) {\textbf{Tóxico-Free} \quad no high-risk};
  \node[badge, below=of b2] (b3) {\textbf{Vegan} \quad $S \geq 80$};
  \node[badge, below=of b3] (b4) {\textbf{Cruelty-Free} \quad $S \geq 80$};
  \node[badge, below=of b4, fill=red!10, draw=ecored]
    (b5) {\textbf{REJECTED} \quad $S < 70$};
\end{tikzpicture}
\caption{Ecological badge assignment thresholds based on Eco-Score.}
\label{fig:badges}
\end{figure}

\subsection{Immutability and Certificates}

Each audit record stored in MongoDB carries a SHA-256 hash:
\begin{lstlisting}[language=Python, caption={Audit hash generation.}]
raw = f"{product_id}-{provider_id}-
        {eco_score}-{timestamp}"
audit_hash = hashlib.sha256(
               raw.encode()).hexdigest()
\end{lstlisting}
Any post-hoc modification of the audit data invalidates the hash,
providing tamper detection. Approved products can download a
multi-page PDF certificate (generated with ReportLab) containing
the audit hash, ingredient table, risk analysis, and a verification
section, as shown in Figure~\ref{fig:auditflow}.

% ─── TikZ audit workflow ──────────────────────────────────────────
\begin{figure}[H]
\centering
\begin{tikzpicture}[
  node distance=0.45cm,
  font=\scriptsize,
  step/.style={
    rectangle, rounded corners=3pt, draw, thick,
    minimum width=3.5cm, minimum height=0.55cm,
    fill=ecobg, draw=ecogreen, text centered
  },
  decision/.style={
    diamond, draw, thick, fill=yellow!15, draw=ecoyellow!80,
    minimum width=2.2cm, minimum height=0.6cm,
    inner sep=0pt, text centered, font=\tiny
  },
  arr/.style={-{Latex[length=2.5mm]}, thick, color=ecogray}
]
  \node[step] (a) {Producer submits product};
  \node[step, below=of a] (b) {Ingredient list parsed};
  \node[step, below=of b] (c) {Chemical DB scan};
  \node[decision, below=0.5cm of c] (d) {$S \geq 70$?};
  \node[step, below left=0.55cm and 0.8cm of d,
        fill=ecolightgreen!70, draw=ecogreen] (e)
    {APPROVED\\Badges assigned};
  \node[step, below right=0.55cm and 0.8cm of d,
        fill=red!10, draw=ecored] (f)
    {REJECTED\\Observaciones};
  \node[step, below=0.55cm of e, minimum width=3.5cm] (g)
    {SHA-256 hash stored\\in MongoDB};
  \node[step, below=of g] (h)
    {PDF Certificate generated};

  \draw[arr] (a)--(b);
  \draw[arr] (b)--(c);
  \draw[arr] (c)--(d);
  \draw[arr] (d) -- node[left,  font=\tiny]{Yes} (e);
  \draw[arr] (d) -- node[right, font=\tiny]{No}  (f);
  \draw[arr] (e)--(g);
  \draw[arr] (g)--(h);
\end{tikzpicture}
\caption{Audit pipeline from product submission to certificate generation.}
\label{fig:auditflow}
\end{figure}

% ───────────────────────────────────────────────────────────────────
\section{Database Architecture}
% ───────────────────────────────────────────────────────────────────

\subsection{PostgreSQL — Relational Layer}

PostgreSQL hosts the core relational data across six tables
(Figure~\ref{fig:erd}).
Notable design decisions: \texttt{products.eco\_score} carries a
\texttt{CHECK(0--100)} constraint; \texttt{payments.idempotency\_key}
(UUID UNIQUE) prevents duplicate charges; and
\texttt{payment\_splits} enforces an 85/15 revenue model at the
database level.

% ─── TikZ ERD ─────────────────────────────────────────────────────
\begin{figure*}[ht]
\centering
\begin{tikzpicture}[
  font=\scriptsize,
  entity/.style={
    rectangle, draw, thick, rounded corners=3pt,
    fill=blue!6, draw=codeblue,
    minimum width=3.0cm,
    text width=2.8cm,
    inner sep=4pt
  },
  rel/.style={-{Latex[length=2mm]}, thick, color=ecogray}
]

\node[entity] (users) at (0,0) {
  \textbf{\color{codeblue}users}\\
  \rule{\linewidth}{0.4pt}
  id \textit{(UUID PK)}\\
  email UNIQUE\\
  full\_name, phone\\
  role, eco\_score\\
  provider, google\_id
};

\node[entity] (providers) at (4.5,0) {
  \textbf{\color{codeblue}providers}\\
  \rule{\linewidth}{0.4pt}
  id \textit{(UUID PK)}\\
  user\_id \textit{(FK)}\\
  business\_name\\
  ruc UNIQUE\\
  certification JSONB
};

\node[entity] (products) at (9,0) {
  \textbf{\color{codeblue}products}\\
  \rule{\linewidth}{0.4pt}
  id \textit{(UUID PK)}\\
  provider\_id \textit{(FK)}\\
  name, price, stock\\
  eco\_score CHECK~0-100\\
  status, images[]
};

\node[entity] (orders) at (0,-3.8) {
  \textbf{\color{codeblue}orders}\\
  \rule{\linewidth}{0.4pt}
  id \textit{(UUID PK)}\\
  customer\_id \textit{(FK)}\\
  status, total\_amount\\
  platform\_fee (15\%)\\
  paid\_at
};

\node[entity] (order_items) at (4.5,-3.8) {
  \textbf{\color{codeblue}order\_items}\\
  \rule{\linewidth}{0.4pt}
  id \textit{(UUID PK)}\\
  order\_id \textit{(FK)}\\
  product\_id \textit{(FK)}\\
  quantity, unit\_price\\
  subtotal
};

\node[entity] (payments) at (9,-3.8) {
  \textbf{\color{codeblue}payments}\\
  \rule{\linewidth}{0.4pt}
  id \textit{(UUID PK)}\\
  order\_id \textit{(FK)}\\
  method (yape|plin...)\\
  idempotency\_key UNIQUE\\
  amount, status
};

\draw[rel] (providers.west) -- (users.east)
  node[midway, above, font=\tiny]{1:1};
\draw[rel] (products.north west) -- (providers.north east)
  node[midway, above, font=\tiny]{N:1};
\draw[rel] (orders.north) -- (users.south)
  node[midway, left, font=\tiny]{N:1};
\draw[rel] (order_items.west) -- (orders.east)
  node[midway, above, font=\tiny]{N:1};
\draw[rel] (order_items.north) -- (products.south)
  node[midway, right, font=\tiny]{N:1};
\draw[rel] (payments.west) -- (orders.east)
  node[midway, above, font=\tiny]{1:1};

\end{tikzpicture}
\caption{Entity-relationship diagram for the EcoMarket PostgreSQL schema.
  Arrows indicate foreign-key direction (many-to-one or one-to-one).}
\label{fig:erd}
\end{figure*}

\subsection{MongoDB — Audit Ledger}

The \texttt{ecomarket\_audit} database contains a single
\texttt{audits} collection. Each document is keyed by
\texttt{product\_id} and stores the full audit result including
\texttt{eco\_score}, \texttt{badges}, \texttt{audit\_hash},
\texttt{status}, and a nested \texttt{details} object with
per-ingredient analysis and chemical findings. Documents are
upserted on re-audit, but the original hash is preserved
in the history array.

% ───────────────────────────────────────────────────────────────────
\section{Payment Integration}
% ───────────────────────────────────────────────────────────────────

Figure~\ref{fig:payment} illustrates the two payment paths.
The \emph{Stripe path} follows the standard Checkout Session
redirect model and handles webhook signature verification.
The \emph{local path} covers Yape, Plin, and TuPay with
simulated QR and OTP flows suitable for demo and future
real-gateway integration.

\begin{figure}[H]
\centering
\begin{tikzpicture}[
  font=\scriptsize,
  node distance=0.38cm and 0.2cm,
  box/.style={rectangle, rounded corners=3pt, draw, thick,
              minimum width=2.8cm, minimum height=0.5cm,
              text centered, fill=white},
  arr/.style={-{Latex[length=2mm]}, thick, color=ecogray}
]
% Stripe path (left)
\node[box, fill=blue!8, draw=codeblue] (ss) {Create Stripe Session};
\node[box, below=of ss] (sc) {Consumer redirected\\to stripe.com};
\node[box, below=of sc] (sw) {Webhook: \\checkout.session.completed};
\node[box, below=of sw, fill=ecolightgreen!60, draw=ecogreen] (sconf)
  {Order Confirmed};

% Local path (right)
\node[box, right=0.6cm of ss, fill=orange!10, draw=orange!70] (ls)
  {Select Yape / Plin};
\node[box, below=of ls] (lq) {QR / Phone\\displayed};
\node[box, below=of lq] (lp) {POST\\process-local};
\node[box, below=of lp, fill=ecolightgreen!60, draw=ecogreen] (lconf)
  {Order Confirmed};

\draw[arr] (ss)--(sc);
\draw[arr] (sc)--(sw);
\draw[arr] (sw)--(sconf);
\draw[arr] (ls)--(lq);
\draw[arr] (lq)--(lp);
\draw[arr] (lp)--(lconf);

\node[above=0.1cm of ss, font=\scriptsize\bfseries, color=codeblue]
  {Stripe (Card)};
\node[above=0.1cm of ls, font=\scriptsize\bfseries, color=orange!80]
  {Local (Yape/Plin)};
\end{tikzpicture}
\caption{Dual payment flow: Stripe Checkout Session redirect
  (left) and local Peruvian payment simulation (right).}
\label{fig:payment}
\end{figure}

Revenue is split at payment confirmation: providers receive
85\,\% of the product subtotal; the platform retains 15\,\%,
recorded in the \texttt{payment\_splits} table for future
disbursement accounting.

% ───────────────────────────────────────────────────────────────────
\section{Deployment Architecture}
% ───────────────────────────────────────────────────────────────────

\subsection{Production Deployment}

The system uses a hybrid cloud deployment strategy
(Figure~\ref{fig:deploy}):

\begin{itemize}[leftmargin=1.2em, itemsep=2pt]
  \item \textbf{Frontend:} Vercel (CDN edge network, automatic
    CI/CD from \texttt{main} branch).
  \item \textbf{Backend services:} Render.com, each service
    deployed from the multi-stage Dockerfile.
  \item \textbf{Databases:} Neon PostgreSQL (serverless, region
    \texttt{us-east-1}) and MongoDB Atlas (shared cluster).
  \item \textbf{Redis:} Upstash serverless Redis for rate-limit
    counters and session caching.
\end{itemize}

\begin{figure}[H]
\centering
\begin{tikzpicture}[
  font=\scriptsize,
  cloud/.style={ellipse, draw, thick, minimum width=2.2cm,
                minimum height=0.7cm, text centered},
  arr/.style={-{Latex[length=2mm]}, thick, color=ecogray}
]
  \node[cloud, fill=blue!10,  draw=codeblue]  (vercel)  {Vercel\\Frontend};
  \node[cloud, fill=green!10, draw=ecogreen, right=1.2cm of vercel]
    (render) {Render\\Backend};
  \node[cloud, fill=yellow!10, draw=ecoyellow!80,
        below left=1.0cm and 0.2cm of render]
    (neon) {Neon\\PostgreSQL};
  \node[cloud, fill=green!8, draw=ecogreen!70,
        below right=1.0cm and 0.2cm of render]
    (atlas) {MongoDB\\Atlas};
  \node[cloud, fill=red!8, draw=ecored!70,
        below=1.0cm of render]
    (redis) {Upstash\\Redis};

  \draw[arr] (vercel) -- node[above, font=\tiny]{HTTPS} (render);
  \draw[arr] (render) -- (neon);
  \draw[arr] (render) -- (atlas);
  \draw[arr] (render) -- (redis);
\end{tikzpicture}
\caption{Production cloud deployment topology.}
\label{fig:deploy}
\end{figure}

\subsection{Local Development with Docker Compose}

All 11 services (including PostgreSQL, MongoDB, and Redis)
are orchestrated with a single \texttt{docker-compose.yml}.
The multi-stage Dockerfile compiles Java JARs (Maven),
installs Node.js dependencies, and installs Python requirements,
producing a monolithic image managed by Supervisor for
all-in-one Render deployments.

% ───────────────────────────────────────────────────────────────────
\section{Software Functionalities}
% ───────────────────────────────────────────────────────────────────

\subsection{Role-Based Workflows}

EcoMarket supports three roles, each with a distinct workflow:

\begin{enumerate}[leftmargin=1.5em, itemsep=3pt]
  \item \textbf{Producer:} Multi-step registration (personal +
    corporate data, RUC validation); product catalog management
    (name, price, stock, images, ingredient list); real-time
    audit status tracking; downloadable PDF certificates.
  \item \textbf{Consumer:} Browse catalog with Eco-Score filters
    and ecological badges; view per-ingredient safety analysis;
    add to cart; checkout with Yape/Plin/Stripe; track orders.
  \item \textbf{Administrator:} Review pending products; approve
    or reject with written observations; view immutable audit
    ledger; access platform statistics (total sales, active
    clients, eco-score averages).
\end{enumerate}

\subsection{Security}

\begin{itemize}[leftmargin=1.2em, itemsep=2pt]
  \item JWT HS256, 24-hour expiration, shared secret across services.
  \item BCrypt password hashing (cost factor 12).
  \item Per-IP rate limiting via Kong plugins (30/min auth;
    15/min payments).
  \item Stripe webhook HMAC-SHA256 signature verification.
  \item Content Security Policy headers on all Next.js routes.
  \item Input sanitization (SQL injection prevention in Spring Boot;
    DOMPurify XSS prevention in frontend).
  \item Cloudflare Turnstile CAPTCHA on registration and login.
\end{itemize}

% ───────────────────────────────────────────────────────────────────
\section{Required Metadata}
% ───────────────────────────────────────────────────────────────────

\begin{table}[H]
\centering
\caption{Code metadata (SoftwareX format).}
\renewcommand{\arraystretch}{1.15}
\scriptsize
\begin{tabularx}{\columnwidth}{cl X}
\toprule
\textbf{Nr.} & \textbf{Description} & \textbf{Value} \\
\midrule
C1 & Code version   & v1.0.0 \\
C2 & Repository     & \href{https://github.com/JosepVl123463/EcoMarketFinalSoftware}
                        {\texttt{github.com/JosepVl123463/\\EcoMarketFinalSoftware}} \\
C3 & License        & MIT \\
C4 & Version system & Git \\
C5 & Languages      & TypeScript, Java, Python, JavaScript \\
C6 & Frameworks     & Next.js, Spring Boot, FastAPI, Express \\
C7 & Databases      & PostgreSQL, MongoDB, Redis \\
C8 & DevOps         & Docker, Vercel, Render, GitHub Actions \\
C9 & Support email  & \href{mailto:josep.garate@est.unap.edu.pe}
                        {\texttt{josep.garate@est.unap.edu.pe}} \\
\bottomrule
\end{tabularx}
\end{table}

\begin{table}[H]
\centering
\caption{Software metadata (SoftwareX format).}
\renewcommand{\arraystretch}{1.15}
\scriptsize
\begin{tabularx}{\columnwidth}{cl X}
\toprule
\textbf{Nr.} & \textbf{Description} & \textbf{Value} \\
\midrule
S1 & Software version & v1.0.0 \\
S2 & Live URL & \href{https://eco-market-final-software-s2w6.vercel.app}
                 {\texttt{eco-market-final-software-s2w6.vercel.app}} \\
S3 & License  & MIT \\
S4 & Platform & Modern web browser (Chrome, Firefox, Safari, Edge) \\
S5 & End-user deps & None (server-side rendering) \\
S6 & Developer deps & Docker, Node.js 20+, Java 21, Python 3.11+ \\
S7 & Support email & \href{mailto:josep.garate@est.unap.edu.pe}
                      {\texttt{josep.garate@est.unap.edu.pe}} \\
\bottomrule
\end{tabularx}
\end{table}

% ───────────────────────────────────────────────────────────────────
\section{Illustrative Examples}
% ───────────────────────────────────────────────────────────────────

\subsection{Producer Scenario}

A quinoa producer from Puno registers with their RUC, uploads a product
photo, fills in the ingredient list (e.g., ``quinoa orgánica, sal marina''),
and submits for audit. The audit engine scans the ingredients, finds no
prohibited chemicals, calculates $S = 100$, assigns
\textit{Eco-Friendly} and \textit{Tóxico-Free} badges, stores the
SHA-256 hash, and marks the product APPROVED. The producer can then
download the PDF certificate to display on their packaging.

\subsection{Consumer Scenario}

A consumer browses the catalog, filters by $S \geq 90$, and finds the
approved quinoa. They view its per-ingredient safety table, add it to
the cart, and proceed to checkout. They select Yape, scan the QR code
on their mobile, and enter the 6-digit approval code. The platform
confirms the order, decrements stock in product-service, and sends
a push notification: \textit{``¡Pago Exitoso! Tu orden \#XYZ ha sido
procesada.''} A 15\,\% platform fee is recorded in
\texttt{payment\_splits}.

% ───────────────────────────────────────────────────────────────────
\section{Impact and Conclusions}
% ───────────────────────────────────────────────────────────────────

EcoMarket demonstrates the feasibility of combining rigorous
ecological auditing with practical digital commerce for underserved
local markets. The platform:

\begin{itemize}[leftmargin=1.2em, itemsep=2pt]
  \item Removes intermediaries for Puno's ecological producers,
    enabling 85\,\% revenue retention.
  \item Provides consumers with verifiable, cryptographically
    anchored product quality information.
  \item Establishes a technical template — microservices, AI
    auditing, local payments — replicable for other
    Peruvian ecological regions.
\end{itemize}

Future work will focus on: (i) integration of a real
computer-vision ingredient detection model; (ii) native mobile
application; (iii) integration with additional Peruvian payment
processors (Culqi, Niubiz); and (iv) stronger certificate
verification via blockchain anchoring.

% ───────────────────────────────────────────────────────────────────
\section*{Acknowledgements}
% ───────────────────────────────────────────────────────────────────

The author thanks the Universidad Nacional del Altiplano, Puno, and the
Escuela Profesional de Ingeniería de Sistemas for their academic support.
Special gratitude is expressed to Professor Torres Cruz Fred, instructor
of the Taller de Software course, for his guidance throughout the
development of EcoMarket. The author also thanks local ecological
producers and early users who provided feedback to improve the platform.

% ───────────────────────────────────────────────────────────────────
\begin{thebibliography}{9}
\bibitem{nextjs} Vercel Inc., \textit{Next.js 16 Documentation},
  \url{https://nextjs.org/docs}, 2026.
\bibitem{springboot} VMware Inc., \textit{Spring Boot 3.2 Reference Guide},
  \url{https://docs.spring.io/spring-boot/docs/3.2}, 2024.
\bibitem{fastapi} S. Ramírez, \textit{FastAPI Documentation},
  \url{https://fastapi.tiangolo.com}, 2024.
\bibitem{kong} Kong Inc., \textit{Kong Gateway 3.6 Documentation},
  \url{https://docs.konghq.com/gateway/3.6}, 2024.
\bibitem{stripe} Stripe Inc., \textit{Stripe Checkout Documentation},
  \url{https://stripe.com/docs/checkout}, 2024.
\bibitem{reportlab} ReportLab Inc., \textit{ReportLab PDF Library},
  \url{https://www.reportlab.com/docs/reportlab-userguide.pdf}, 2024.
\bibitem{cosmos} COSMOS, \textit{COSMOS Organic Standard v3.0},
  COSMOS-standard AISBL, Brussels, 2023.
\bibitem{cereach} European Commission, \textit{REACH Regulation
  (EC) No 1907/2006}, Official Journal of the EU, 2006.
\bibitem{docker} Docker Inc., \textit{Docker Compose v3 Reference},
  \url{https://docs.docker.com/compose}, 2024.
\end{thebibliography}

\end{document}
